% ---
% id: meat-core/rem:target-conditions
% kind: rem
% title: "The two index conditions of the removal checklist"
% summary: >-
%   Without (W2) every state-blind constraint-active sector witnesses its own
%   removal, and deleting it can overstate the value by a bit per round.
%   Distinct targets cost no removal in exact arithmetic: repeated targets
%   merge into one witness certifying the same identities.
% prominence: standard
% status: accepted
% uses: [def:classes, thm:removal, prop:retention]
% mentions: [rem:classes]
% assumes: []
% refs: []
% provenance: newly-derived
% lean:
% ---

Condition (W2) of Definition~\ref{def:classes} asks two things of a witness
for $s^{*}$: that it does not name $s^{*}$, and that it names no sector
twice. Both come from the construction in the proof of
Theorem~\ref{thm:removal}.

\begin{remark}\label{rem:target-conditions}
\emph{Without (W2), every state-blind constraint-active sector witnesses
its own removal.} Let $s^{*}$ be state-blind and key-inert for every
$\hat c\in\Chkey$ --- that is, state-blind and constraint-active in the
sense of Definition~\ref{def:classes}\,(b) --- and take the
\emph{self-witness} $L=1$, $s_1\coloneqq s^{*}$, $\mu_1\coloneqq1$, and
$\chi_1$ any state on $\Qsys_{s^{*}}$ in the form (W4) requires --- a
diagonal projector $\dyad{i}$ will do. Then (W1) holds by hypothesis, (W3)
and (W4) hold by construction, (W5) holds because $s^{*}$ is itself
key-inert for every $\hat c\in\Chkey$, and \emph{both equality families
close exactly}: state-blindness gives
$\tr[F^{(s^{*})}_v\chi_1]=c_{v,s^{*}}\tr\chi_1=c_{v,s^{*}}$ and
$\tr[\hat O^{(s^{*})}_k\chi_1]=o_{k,s^{*}}\tr\chi_1=o_{k,s^{*}}$, and the
evaluation is the product $c_{v,s^{*}}\cdot1$, which rounds nothing, so
each residual is the datum $0$ and (W6) passes as well. Every condition
but (W2) is met, and a checklist without (W2) would flag the sector
removable and delete a constraint-active sector --- the one operation
Proposition~\ref{prop:retention} proves unsound. Its counterexample is of
exactly this shape: its sector~2 has $F^{(2)}_0=1$, $F^{(2)}_1=0$ on a
one-dimensional block, hence is state-blind, and is key-inert, hence
constraint-active; deleting it \emph{overstates $\hlow$ by one bit per
round}. What forbids the self-witness is the construction in the proof of
Theorem~\ref{thm:removal}: it sets $\hat\omega_{s^{*}}\coloneqq0$ and
$\hat\omega_{s_i}\coloneqq\omega_{s_i}+\mu_i\,m\otimes\chi_i$
\emph{simultaneously}, and if some $s_i=s^{*}$ those two assignments
contradict each other, so $\hat\omega$ is not defined and no step of the
proof is available. The theorem quantifies over
$s_1,\dots,s_L\in\SecSet\setminus\{s^{*}\}$ for that reason, and the
first half of (W2) is that quantifier and nothing else.

\emph{Distinct targets, at no cost.} With $s_i=s_j$ for $i\ne j$ the
proof's second assignment
$\hat\omega_{s_i}\coloneqq\omega_{s_i}+\mu_i\,m\otimes\chi_i$ would assign
two values to one block, the same collision the self-witness produces at
$s_i=s^{*}$, and $\hat\omega$ would again not be defined. The theorem
therefore quantifies over pairwise distinct targets and (W2) is that
quantifier as well. \emph{No removal is lost by it.} A candidate witness
with repeated targets \emph{merges}: drop every term with $\mu_i=0$, and
replace the terms at a repeated target $s$ by the single term
$(\bar\mu_s,\bar\chi_s)$ with $\bar\mu_s\coloneqq\sum_{i:s_i=s}\mu_i$ and
$\bar\chi_s\coloneqq\bar\mu_s^{-1}\sum_{i:s_i=s}\mu_i\chi_i$, defined
because $\bar\mu_s>0$ after the drop. Then $\bar\mu_s\ge0$ and
$\sum_s\bar\mu_s=1$, so (W3) holds; $\bar\chi_s\succeq0$ with
$\tr\bar\chi_s=1$, a convex combination of states, so (W4) holds; and
$\bar\mu_s\tr[F^{(s)}_v\bar\chi_s]=\sum_{i:s_i=s}\mu_i\tr[F^{(s)}_v\chi_i]$
by linearity of $X\mapsto\tr[F^{(s)}_vX]$, and likewise for
$\hat O^{(s)}_k$, so (RM1) and (RM3) take the same values and (W6)'s
identities are the same identities. \emph{In exact arithmetic the merged
witness certifies exactly what the original would have.} The merged
witness and its evaluation order must still be \emph{exhibited}, because
(W3), (W4) and (W6) are conditions on the exhibited numbers: a merge whose
arithmetic rounds sets no flag, which is the safe direction
(Remark~\ref{rem:classes}).
\end{remark}
